\documentclass[10pt]{beamer}

% Theme Choice
\usetheme{Madrid} 
\usecolortheme{seahorse} % Built-in, light blue and professional

% Necessary Packages
\usepackage[utf8]{inputenc}
\usepackage{amsmath, amssymb, amsthm} % For heavy math
\usepackage{bookmark, booktabs, array}
\usepackage{adjustbox}
\usepackage{tikz} % For drawing your graphs
\usetikzlibrary{graphs, graphs.standard, arrows.meta}

\AtBeginSection[]
{
  \begin{frame}
    \vfill
    \centering
    \begin{beamercolorbox}[sep=8pt,center,shadow=true,rounded=true]{title}
        \usebeamerfont{title}\insertsectionhead\par%
    \end{beamercolorbox}
    \vfill
  \end{frame}
}

% Custom Block for Your "Board Work" transition
\setbeamercolor{block title}{bg=blue!70,fg=white}
\setbeamercolor{block body}{bg=blue!10,fg=black}

\title[Homomorphism Counting]{Counting Graph Homomorphisms in Bipartite Settings}
\subtitle{Paper Presentation: CS1.502 - Information Theoretic Methods in Computer Science (2026)}
\author{Shreyas Kasture  \\ Harshdeep Singh }
\institute{IIIT Hyderabad}
\date{\today}

\begin{document}

\begin{frame}
  \titlepage
\end{frame}


% ============================================================
\section{Introduction \& Setup}
% ============================================================

\begin{frame}{Problem Setting}
  \begin{block}{Setup}
    \begin{itemize}
      \item \textbf{Source graph:} $K_{p,q}$ (complete bipartite)
      \item \textbf{Target graph:} $G$ bipartite with partite sets $L$, $R$ of sizes $n_1$, $n_2$
      \item \textbf{Goal:} Bound $\hom(K_{p,q}, G)$ --- the number of graph homomorphisms
    \end{itemize}
  \end{block}

  \begin{block}{Notation}
    \begin{itemize}
      \item $\delta = |E(G)|/(n_1 n_2)$: edge density of $G$
      \item $d_G(w)$: degree of vertex $w$ in $G$
      \item $N_{k,\ell}(G)$: number of labeled bipartite cliques $K_{k,\ell}$ in $G$ with sizes $k$, $\ell$
    \end{itemize}
  \end{block}

  \begin{alertblock}{Challenge}
    Direct enumeration is computationally expensive --- we need efficient bounds.
  \end{alertblock}
\end{frame}

\begin{frame}{Two Complementary Approaches}
  \begin{columns}
    \begin{column}{0.48\textwidth}
      \begin{block}{Section 3: Combinatorial Bounds}
        \begin{itemize}
          \item Exact formula via Stirling numbers
          \item Simplified lower bound
          \item Uses graph structure directly
          \item Exact on $C_4$-free graphs
        \end{itemize}
      \end{block}
    \end{column}
    \begin{column}{0.48\textwidth}
      \begin{alertblock}{Section 4: Entropy-Based Bounds}
        \begin{itemize}
          \item Probabilistic edge sampling
          \item Shannon entropy machinery
          \item Simple bound from density $\delta$ only
          \item Refined bound using degree distribution
        \end{itemize}
      \end{alertblock}
    \end{column}
  \end{columns}
  \vspace{1em}
  \textbf{Key insight:} More structural information $\Rightarrow$ tighter bounds, but higher cost.
\end{frame}

% ============================================================
\section{Section 3:Combinatorial Bounds}
% ============================================================

\begin{frame}{Proposition 3.1: Exact Formula}
  \begin{theorem}[Proposition 3.1]
    Let $G$ be bipartite with partite sets $L$, $R$. For $p, q \in \mathbb{N}$:
    \[
      \hom(K_{p,q}, G) = \sum_{k=1}^{p} \sum_{\ell=1}^{q} k!\, \ell!\, S(p,k)\, S(q,\ell)
      \bigl[N_{k,\ell}(G) + N_{\ell,k}(G)\bigr]
    \]
  \end{theorem}

  \textbf{Proof outline:}
  \begin{enumerate}
    \item \textbf{Bipartition preservation:} Every homomorphism maps each partite set of $K_{p,q}$ entirely into one partite set of $G$.
    \item \textbf{Image is a clique:} The image of $\varphi$ forms a complete bipartite subgraph $K_{k,\ell} \subseteq G$.
    \item \textbf{Counting surjections:} For each $K_{k,\ell}$ in $G$, there are $k!\,S(p,k) \cdot \ell!\,S(q,\ell)$ surjective mappings.
    \item \textbf{Sum over both orientations} ($N_{k,\ell}$ and $N_{\ell,k}$), then sum over all $k, \ell$.
  \end{enumerate}
\end{frame}

\begin{frame}{Proposition 3.2: Combinatorial Lower Bound}
  \begin{theorem}[Proposition 3.2]
    \[
      \hom(K_{p,q}, G) \geq \sum_{w \in V(G)} \bigl[d_G(w)^p + d_G(w)^q\bigr] - 2|E(G)|
    \]
    with equality if and only if $G$ is $C_4$-free.
  \end{theorem}

  \textbf{Proof strategy:}
  \begin{enumerate}
    \item \textbf{Drop higher terms} in Prop.\ 3.1: all $k, \ell \geq 2$ terms are non-negative, so dropping them gives a lower bound.
    \item \textbf{Star homomorphisms:} The remaining sum equals
      $\hom(K_{1,p}, G) + \hom(K_{1,q}, G) - 2|E(G)|$.
    \item \textbf{Classical result:} $\hom(K_{1,m}, G) = \sum_w d_G(w)^m$.
    \item \textbf{Equality:} Holds $\Leftrightarrow$ all dropped terms vanish $\Leftrightarrow$ $N_{2,2}(G) = 0$ $\Leftrightarrow$ $G$ is $C_4$-free.
  \end{enumerate}
\end{frame}

\begin{frame}{Remarks 3.2 \& 3.3: Sparsity Requirement for Exactness}
  \begin{block}{Remark 3.2 --- Edge Count in $C_4$-Free Graphs}
    By the K\H{o}v\'ari--S\'os--Tur\'an theorem, any $C_4$-free bipartite graph satisfies
    \[
      m = O(n^{3/2})
    \]
    So the equality condition in Prop.\ 3.2 forces $G$ to be very sparse.
  \end{block}

  \begin{block}{Remark 3.3 --- Edge Density Constraint}
    Translating to density: $\delta = m/(n_1 n_2) = O(n^{-1/2}) \to 0$ as $n \to \infty$.\\[0.3em]
    \textbf{Consequence:} The combinatorial bound is exact only for graphs that become \emph{increasingly sparse} as they grow.
  \end{block}

  \begin{alertblock}{Takeaway}
    Prop.\ 3.2 is well-suited for sparse, $C_4$-free graphs. For dense graphs, we need a different approach.
  \end{alertblock}
\end{frame}

\begin{frame}{Remark 3.4: Failure on Dense Graphs}
  \textbf{Test case:} $G = K_{n/2,\,n/2}$ (both partite sets of size $n/2$, all edges present).

  \begin{columns}
    \begin{column}{0.52\textwidth}
      \textbf{Exact count:}
      \[
        \hom(K_{p,p},\,G) = 2^{1-p}\,n^{2p}
      \]
      Each of $p$ left-vertices maps independently to $n/2$ choices on each side.

      \textbf{Prop.\ 3.2 lower bound:}
      \[
        \hom(K_{p,p},\,G) \gtrsim 2^{1-p}\,n^{p+1}
      \]
    \end{column}
    \begin{column}{0.44\textwidth}
      \begin{alertblock}{Ratio}
        \[
          \frac{\text{Exact}}{\text{Bound}} \;\sim\; \frac{n^{2p}}{n^{p+1}} = n^{p-1} \;\to\; \infty
        \]
        The bound becomes \textbf{arbitrarily weak}.
      \end{alertblock}
    \end{column}
  \end{columns}

  \vspace{0.5em}
  \textit{``This is in contrast to the lower bounds derived in Section~4, for which the corresponding ratio tends, as desired, to~1.''}
\end{frame}

\begin{frame}{Corollaries 3.1 \& 3.2}
  \begin{block}{Corollary 3.1 --- Trees (equality case)}
    For any tree $T$ on $n$ vertices (always $C_4$-free):
    \[
      \hom(K_{p,q},\,T) = \sum_{w \in V(T)} \bigl[d_T(w)^p + d_T(w)^q\bigr] - 2(n-1)
    \]
  \end{block}

  \begin{block}{Corollary 3.2 --- Jensen's Inequality Refinement}
    Since $x \mapsto x^p$ is convex, Jensen's inequality applied to vertex degrees gives:
    \[
      \hom(K_{p,q},\,G) \;\geq\; 2^p\,n^{1-p}\,m^p + 2^q\,n^{1-q}\,m^q - 2m
    \]
    \textbf{Equality} iff $G$ is $C_4$-free \emph{and} regular (all degrees equal, so Jensen is tight).
  \end{block}
\end{frame}

\begin{frame}{Corollaries 3.3 \& 3.4: Cycles and Paths}
  \begin{block}{Corollary 3.3 --- Cycle $C_n$}
    \[
      \hom(K_{p,q},\,C_n) = \begin{cases}
        n\bigl(2^p + 2^q - 2\bigr) & n = 3 \text{ or } n \geq 5 \\
        2^{p+q+1} & n = 4
      \end{cases}
    \]
    \textit{Why does this work for non-bipartite odd cycles?}
    A bipartite source $K_{p,q}$ cannot create odd cycles, so its image lies in a tree-like (path) subgraph of $C_n$, which is $C_4$-free --- the Prop.\ 3.2 formula applies automatically.
    $C_4$ is handled by direct counting ($N_{k,\ell}(C_4) = 0$ for $k,\ell \geq 2$).
  \end{block}

  \begin{block}{Corollary 3.4 --- Path $P_\ell$}
    $P_\ell$ is a tree ($\ell - 2$ vertices of degree~2, two of degree~1), so Corollary~3.1 applies directly:
    \[
      \hom(K_{p,q},\,P_\ell) = (\ell-2)(2^p+2^q) + 2 - 2(\ell-1)
    \]
  \end{block}
\end{frame}

\begin{frame}{Transition: Why Entropy-Based Bounds?}
  \begin{itemize}
    \item Prop.\ 3.2 requires knowledge of the full degree sequence --- $O(n)$ input.
    \item Prop.\ 3.1 requires the full biclique structure --- $O(n^{k+\ell})$ computation.
    \item \textbf{Both fail on dense graphs:} ratio of exact count to bound $\to \infty$.
  \end{itemize}

  \vspace{0.5em}
  \begin{alertblock}{Section 4 Goal}
    Derive lower bounds using \textbf{only edge density $\delta$} (and optionally the degree distribution) --- no structural enumeration. The resulting ratio approaches \textbf{1} on balanced, dense graphs.
  \end{alertblock}

  \vspace{0.3em}
  \textbf{Tool:} Shannon entropy applied to a random edge $(U, V) \sim \mathrm{Uniform}(E(G))$.
\end{frame}

% ============================================================
\section{Section 4: Entropy-Based Bounds}
% ============================================================

\begin{frame}{Entropy Background}
  \begin{block}{Definition (Shannon Entropy)}
    For discrete r.v.\ $X$ with support $[n]$:
    \[
      H(X) = -\sum_{i=1}^{n} P(X=i) \log P(X=i)
    \]
  \end{block}

  \textbf{Key properties used:}
  \begin{itemize}
    \item $H(X) \leq \log n$ (uniform distribution maximizes entropy)
    \item $H(X,Y) = H(X) + H(Y \mid X)$ (chain rule)
    \item $H(X_1, \ldots, X_n) \leq \sum_i H(X_i)$ (subadditivity)
  \end{itemize}

  \textbf{Strategy:} Sample edges randomly from $G$, build a joint distribution $(U^p, V^q)$, and use entropy lower bounds to count homomorphisms.
\end{frame}

\begin{frame}{Lemma 4.1: Probabilistic Construction}
  \begin{block}{Setup}
    Let $(U,V)$ be uniform over $E(G)$, $U \in L$, $V \in R$.
  \end{block}

  \textbf{Construct $(U^p, V^q)$ as follows:}
  \begin{enumerate}
    \item $V_1, \ldots, V_q$ are \textbf{i.i.d.\ given} $U$: \quad
      $P_{V^q \mid U}(v \mid u) = \prod_j P_{V \mid U}(v_j \mid u)$
    \item $U_1, \ldots, U_p$ are \textbf{i.i.d.\ given} $V^q$ via Bayes-like construction.
  \end{enumerate}

  \begin{theorem}[Lemma 4.1]
    \textit{Under this construction:}\\[0.3em]
    \quad (i) $U_i \sim U$ for all $i \in [p]$\\
    \quad (ii) $(U_i, V_j) \sim (U,V)$ for all $i, j$
  \end{theorem}

  \textbf{Importance:} Since $(U_i, V_j)$ has the correct edge distribution, any realization
  $(u,v) \in U^p \times V^q$ yields a valid homomorphism $\varphi : V(K_{p,q}) \to V(G)$.
\end{frame}

\begin{frame}{Lemma 4.2: Entropy Lower Bounds}
  \begin{theorem}[Lemma 4.2]
    \[
      H(U_1, V^q) \geq \log(\delta^q n_1 n_2^q)
    \]
    \[
      H(U^p, V^q) \geq \log(\delta^{pq} n_1^p n_2^q)
    \]
  \end{theorem}

  \textbf{Proof of first inequality:}
  \begin{align*}
    H(U_1, V^q) &= H(U) + q\,H(V \mid U) \quad \text{(chain rule + cond.\ independence)}\\
                &= H(U) + q\bigl[\log(\delta n_1 n_2) - H(U)\bigr]\\
                &= q\log(\delta n_1 n_2) - (q-1)\,H(U)\\
                &\geq q\log(\delta n_1 n_2) - (q-1)\log n_1 \quad \text{(since $H(U) \leq \log n_1$)}\\
                &= \log(\delta^q n_1 n_2^q)
  \end{align*}

  \textbf{Second inequality:} parallel argument using $H(V^q) \leq q \log n_2$.
\end{frame}

\begin{frame}{Proposition 4.1: First Entropy-Based Bound}
  \begin{theorem}[Proposition 4.1]
    Let $G$ be bipartite with $|L| = n_1$, $|R| = n_2$, density $\delta$. Then:
    \[
      \hom(K_{p,q}, G) \geq \delta^{pq}(n_1^p n_2^q + n_1^q n_2^p)
    \]
  \end{theorem}

  \textbf{Proof:}
  \begin{enumerate}
    \item Partition $\mathrm{Hom}(K_{p,q}, G) = H_1 \sqcup H_2$ by which partite sets map where.
    \item Each realization $(u,v)$ of $(U^p, V^q)$ gives a distinct element of $H_1$.
    \item Uniform distribution: $H(U^p, V^q) \leq \log |H_1|$
    \item Lemma 4.2: $\log(\delta^{pq} n_1^p n_2^q) \leq H(U^p, V^q) \leq \log |H_1|$, so $|H_1| \geq \delta^{pq} n_1^p n_2^q$.
    \item By symmetry: $|H_2| \geq \delta^{pq} n_1^q n_2^p$.
    \item Sum: $\hom(K_{p,q}, G) \geq \delta^{pq}(n_1^p n_2^q + n_1^q n_2^p)$.
  \end{enumerate}
\end{frame}

\begin{frame}{Discussion 4.1: Comparison to Sidorenko's Bound}
  \begin{columns}
    \begin{column}{0.48\textwidth}
      \textbf{Sidorenko's Conjecture (classical):}
      \[
        \hom(K_{p,q}, G) \geq (2\delta)^{pq}(n_1 + n_2)^{p+q-2} \cdots
      \]

      \textbf{Our entropy bound (Prop.\ 4.1):}
      \[
        \hom(K_{p,q}, G) \geq \delta^{pq}(n_1^p n_2^q + n_1^q n_2^p)
      \]
    \end{column}
    \begin{column}{0.48\textwidth}
      \begin{alertblock}{Result}
        \textbf{Ratio} $\dfrac{\text{Prop.\ 4.1}}{\text{Sidorenko}} \geq 2^{p-1}$ for $p = q$.\\[0.4em]
        For $p > q$: improvement $\geq 2^{pq-(p+q)}$.\\[0.3em]
        Grows \textbf{exponentially} in $|p - q|$.
      \end{alertblock}
    \end{column}
  \end{columns}
  \vspace{0.5em}
  The entropy approach is provably superior to Sidorenko for $K_{p,q}$ homomorphism counting.
\end{frame}

\begin{frame}{Proposition 4.2: Refined Entropy Bound}
  \begin{block}{Key idea}
    Replace the uniform bound $H(U) \leq \log n_1$ with the \textit{actual} degree entropy.
  \end{block}

  \textbf{Define degree entropies:}
  \[
    x = H(U) = \sum_k \frac{d_k^{(U)}}{|E(G)|} \log \frac{|E(G)|}{d_k^{(U)}}, \qquad
    y = H(V) = \sum_k \frac{d_k^{(V)}}{|E(G)|} \log \frac{|E(G)|}{d_k^{(V)}}
  \]

  \begin{theorem}[Proposition 4.2]
    \begin{align*}
      \hom(K_{p,q}, G) \geq\;
        &\max\!\bigl\{(\delta n_1 n_2)^{pq} e^{-p(q-1)x - q(p-1)y},\;
          (\delta n_1 n_2)^q e^{-(q-1)x},\;
          (\delta n_1 n_2)^p e^{-(p-1)y}\bigr\}\\
        +\;&\max\!\bigl\{(\delta n_1 n_2)^{pq} e^{-q(p-1)x - p(q-1)y},\;
          (\delta n_1 n_2)^p e^{-(p-1)x},\;
          (\delta n_1 n_2)^q e^{-(q-1)y}\bigr\}
    \end{align*}
  \end{theorem}
\end{frame}

\begin{frame}{Proposition 4.2: Proof Outline \& When It Helps}
  \textbf{Proof outline:}
  \begin{enumerate}
    \item Use exact entropy: $H(U_1, V^q) = q\log(\delta n_1 n_2) - (q-1)H(U)$ without loosening $H(U)$.
    \item Apply chain rule to $H(U^p, V^q)$ to derive multiple lower bounds.
    \item Take max of all derived bounds per orientation; sum over both orientations.
  \end{enumerate}

  \textbf{When does the refinement help?}
  \begin{itemize}
    \item \textbf{Regular graphs:} $H(U) = \log n_1$, $H(V) = \log n_2$ --- reverts to Prop.\ 4.1.
    \item \textbf{Star graph} (one hub of degree $n$):
      \[
        H(U) \ll \log n \implies \text{exponentially tighter bound}
      \]
    \item \textbf{Skewed degrees:} Non-uniform distribution $\Rightarrow$ $H$ significantly below $\log n$.
  \end{itemize}
\end{frame}

% ============================================================
\section{Comparison \& Takeaways}
% ============================================================

\begin{frame}{Summary Comparison}
  \begin{adjustbox}{max width=\linewidth, center}
    \begin{tabular}{lcccc}
      \toprule
      & \textbf{Prop.\ 3.1} & \textbf{Prop.\ 3.2} & \textbf{Prop.\ 4.1} & \textbf{Prop.\ 4.2} \\
      \midrule
      Tightness     & Exact          & Good        & Good           & Best \\
      When exact    & Always         & $C_4$-free  & Extremes only  & Regular + special \\
      Complexity    & $O(n^{k+\ell})$& $O(n)$      & $O(1)$         & $O(n)$ \\
      Input needed  & Full structure & Degrees     & $\delta, n_1, n_2$ & Degree dist. \\
      Beats Sidorenko & ---          & Varies      & $\geq 2^{p-1}\times$ & Even better \\
      \bottomrule
    \end{tabular}
  \end{adjustbox}
\end{frame}

\begin{frame}{Key Takeaways}
  \begin{enumerate}
    \item \textbf{Stirling numbers} quantify partition structure --- central to the exact formula.
    \item \textbf{$C_4$-free condition} makes the combinatorial bound exact: absence of $K_{2,2}$ eliminates all higher biclique terms.
    \item \textbf{Entropy framework} translates counting to probability, enabling systematic lower bounds from density alone.
    \item \textbf{Degree entropy} in Prop.\ 4.2 captures non-regularity, giving exponentially tighter bounds for skewed graphs.
    \item \textbf{Progressive improvement:}
      \begin{multline*}
        \underbrace{\text{Comb.\ (exact, expensive)}}_{\text{Prop.\ 3.1}}
        \to \underbrace{\text{Comb.\ lower bound}}_{\text{Prop.\ 3.2}} \\
        \to \underbrace{\text{Simple entropy}}_{\text{Prop.\ 4.1}}
        \to \underbrace{\text{Refined entropy}}_{\text{Prop.\ 4.2}}
      \end{multline*}
  \end{enumerate}
\end{frame}

\end{document}

